% GENERATED from Aim/Common/conventions.json by scripts/build_notation.py.
% Do not edit this file. Edit the JSON and run the script.
\documentclass[11pt]{article}
\usepackage[margin=1in]{geometry}
\usepackage[T1]{fontenc}
\usepackage{lmodern}
\usepackage{amsmath,amssymb}
\usepackage{booktabs,longtable,array}
\renewcommand{\arraystretch}{1.25}
\title{Notation and Conventions}
\author{Adventures in Mathematics}
\date{}
\begin{document}
\maketitle

\noindent This page lists how letters are used in the Lean sources, where the
Lean departs from standard notation, and the conventions the sources follow.
It is generated from
\texttt{Aim/Common/conventions.json}.

\section*{Letters}
\begin{longtable}{@{}>{\raggedright}p{0.20\linewidth}>{\raggedright}p{0.34\linewidth}>{\raggedright\arraybackslash}p{0.38\linewidth}@{}}
\toprule
Letters & Stand for & Note \\
\midrule
\endhead
X, Y & a type & the ambient type of a set, or the target of a function \\
F & a field or an ordered field & also used for a family of sets in the set-theory files, never in the same file \\
R & a complete ordered field & the real numbers, taken abstractly \\
N & a system of natural numbers & any type satisfying the natural-number interface \\
I & an index type &  \\
S, T & a set & sets are uppercase \\
C, D & a collection: a set of sets &  \\
F, G & an indexed family of sets &  \\
x, y, z, a, b, c & an element & elements are lowercase \\
m, n, k & a natural number &  \\
$\varepsilon$ & a positive element of an ordered field &  \\
p & a property of elements & a function into propositions \\
f, g & a function &  \\
h, hab, hx & a hypothesis & names of proofs begin with h \\
\bottomrule
\end{longtable}

\section*{Where the Lean differs from standard notation}
\noindent Standard mathematical notation is used wherever Lean has it and is
not listed here. These are the exceptions.
\begin{longtable}{@{}>{\raggedright}p{0.26\linewidth}>{\raggedright}p{0.24\linewidth}>{\raggedright\arraybackslash}p{0.42\linewidth}@{}}
\toprule
Standard & In Lean & Why \\
\midrule
\endhead
$S^{c}$ & \texttt{Set.compl S} & no complement notation yet \\
$\bigcup_{i} S_i$ & \texttt{Set.iUnion S} & no big-union notation yet \\
$\bigcap_{i} S_i$ & \texttt{Set.iInter S} & no big-intersection notation yet \\
$\bigcup C$ & \texttt{Set.sUnion C} & no notation yet for the union of a collection \\
$X \text{ (as a set)}$ & \texttt{Set.univ} & a type is not a set; the set of all its elements has a name \\
$|x|$ & \texttt{abs x} & no bar notation yet \\
$a = \sup S$ & \texttt{IsLeastUpperBound S a} & the supremum is a relation until a function is defined \\
$0^{-1} \text{ undefined}$ & \texttt{0$^{-1}$} & the inverse is a total function; its value at 0 is unspecified, so inverse lemmas assume a nonzero element \\
$n \in \mathbb{N} \subseteq \mathbb{R}$ & \texttt{natCast n} & a natural number is sent into a field by an explicit map \\
$n + 1$ & \texttt{succ n} & successor is primitive; addition is defined from it \\
$\mathbb{N} = \{1, 2, 3, \dots\}$ & \texttt{0} & the natural numbers here start at 0; Landau's and Feferman's start at 1 \\
\bottomrule
\end{longtable}

\section*{Names}
\noindent A concept whose Lean name is abbreviated also has a readable full
name. The Lean name is the definition; the full name stands for it.
\begin{longtable}{@{}>{\raggedright}p{0.20\linewidth}>{\raggedright}p{0.32\linewidth}>{\raggedright\arraybackslash}p{0.40\linewidth}@{}}
\toprule
Lean name & Readable name & Meaning \\
\midrule
\endhead
\texttt{IsLeast} & \texttt{IsMinimum} & the least element of a set \\
\texttt{IsGreatest} & \texttt{IsMaximum} & the greatest element of a set \\
\texttt{IsLUB} & \texttt{IsLeastUpperBound} & the least upper bound (supremum) of a set \\
\texttt{IsGLB} & \texttt{IsGreatestLowerBound} & the greatest lower bound (infimum) of a set \\
\texttt{BddAbove} & \texttt{BoundedAbove} & a set that has an upper bound \\
\texttt{BddBelow} & \texttt{BoundedBelow} & a set that has a lower bound \\
\bottomrule
\end{longtable}

\section*{Conventions}
\begin{enumerate}
\item A major concept has a readable full name. Where its Lean name is abbreviated, the Lean name is the definition and the full name is an \texttt{abbrev} for it. Lemmas have one name, after the Lean name. \emph{(checked automatically)}
\item Every theorem, lemma, corollary, proposition and example is proved with a \texttt{by} tactic block, so that the proof can be stepped through in the InfoView. \emph{(checked automatically)}
\item A file that only re-exports its folder is named \textless{}Folder\textgreater{}Aggregator.lean and contains imports and a module docstring, no declarations. \emph{(checked automatically)}
\item Mathematics is declared in \texttt{namespace Aim}. A sub-namespace exists only when it is the name of a type and holds that type's own operations. \emph{(checked automatically)}
\begin{itemize}
\item Exception: \texttt{Aim/Common/Common.lean}, syntax macros, declared at the root.
\item Exception: \texttt{Aim/Common/Tactics.lean}, tactic macros.
\item Exception: \texttt{Aim/Logic/Syntax.lean}, to be rewritten under Logic/Propositional; see Aim/Logic/logic-todo.md.
\item Exception: \texttt{Aim/MetaLanguage/Eq.lean}, compatibility names, to be removed.
\item Exception: \texttt{Aim/IdentityAndEquality/Constructions/}, not yet flattened.
\end{itemize}
\item A file under Aim/ imports the interface files it needs, never the root module \texttt{Aim}, which also brings in every construction. \emph{(checked automatically)}
\item Nothing under Aim/ or tests/ imports Mathlib. Mathlib appears only in validation/mathlib. \emph{(checked automatically)}
\item A variable that is a set, a collection or a family of sets is an uppercase letter. An element of a set is a lowercase letter. \emph{(checked automatically)}
\item Subtraction and division on a field are the derived operations. A construction of a number system never defines its own. \emph{(not checked automatically)}
\item A proof cites the lemmas of an interface (membership lemmas for sets, the two equations for \texttt{recurse}) and does not unfold its definitions. \emph{(not checked automatically)}
\item A \texttt{sorry} marks a proof still to be written by the learner. It is never a defect. \emph{(not checked automatically)}
\end{enumerate}
\end{document}
